% GENERATED FILE. Edit canonical lessons, then run npm run generate:books.
% canonical-source-hash: fnv1a64:b9301fb2154129af
% canonical-lessons: PT-C284-escandalo, PT-C284-censura, PT-C284-editorial, PT-C284-cronica, PT-C284-cronista, PT-R284-revisao-277-284

\chapter{Escândalos}
\label{ch:pt-escandalos}
\hlchaptermodality{284}

\begin{chapteropening}
\textbf{By the end of this chapter:} \emph{I can talk about scandal, censorship, editorials, newspaper pieces and columnists.}

Complete the last lesson of chapter 284: I can talk about scandal, censorship, editorials, newspaper pieces and columnists.
\end{chapteropening}

\section[o escândalo]{o escândalo — the scandal}

\label{lesson:PT-C284-escandalo}



\begin{quote}
Before the new one: say the Portuguese for the denial, then the Portuguese for the spokesperson.
\end{quote}

\subsection*{You'll want to know: o escândalo}
\textbf{o escândalo} — \textquotedblleft{}the scandal\textquotedblright{}.

\begin{quote}
O escândalo saiu em todos os jornais.
\end{quote}

\begin{quote}
\emph{The scandal was in all the papers.}
\end{quote}

\emph{Que escândalo!} — how shocking!

\begin{cousinweb}
\emph{escândalo} is a learned borrowing from Latin \emph{scandalum} (EtymDB), from Greek \emph{skándalon} \textquotedblleft{}a trap, a stumbling block\textquotedblright{}.
\end{cousinweb}

\begin{grammarlens}[title={What you've built}]
One more word for telling, explaining and describing at B1.
\end{grammarlens}

\subsection*{Your turn}
\begin{itemize}
\raggedright
  \item \emph{Say it:} \emph{o escândalo}
  \item \emph{Say it:} \emph{o escândalo}, once more
  \item \emph{Say it:} a sentence of your own with \emph{o escândalo}
  \item \emph{Recall:} say the Portuguese for the denial, then the Portuguese for the spokesperson, then say \emph{o escândalo} again
\end{itemize}

\begin{tcolorbox}[breakable,colback=teal!4,colframe=teal!35!black,title={Before you move on}]
What does o escândalo mean? (\textquotedblleft{}The scandal\textquotedblright{}.) Say it once more.
\end{tcolorbox}
\section[a censura]{a censura — censorship; criticism}

\label{lesson:PT-C284-censura}



\begin{quote}
Before the new one: say the Portuguese for the spokesperson, then the Portuguese for the scandal.
\end{quote}

\subsection*{You'll want to know: a censura}
\textbf{a censura} — \textquotedblleft{}censorship; criticism\textquotedblright{}.

\begin{quote}
No passado, a censura cortava as notícias.
\end{quote}

\begin{quote}
\emph{In the past, censorship cut the news.}
\end{quote}

Also blame: \emph{uma censura} is a reproach.

\begin{cousinweb}
\emph{censura} is a learned borrowing from Latin \emph{cēnsūra} \textquotedblleft{}the office of the censor, a judgement\textquotedblright{} (EtymDB), from \emph{cēnsēre} \textquotedblleft{}to assess, to judge\textquotedblright{}.
\end{cousinweb}

\begin{grammarlens}[title={What you've built}]
One more word for telling, explaining and describing at B1.
\end{grammarlens}

\subsection*{Your turn}
\begin{itemize}
\raggedright
  \item \emph{Say it:} \emph{a censura}
  \item \emph{Say it:} \emph{a censura}, once more
  \item \emph{Say it:} a sentence of your own with \emph{a censura}
  \item \emph{Recall:} say the Portuguese for the spokesperson, then the Portuguese for the scandal, then say \emph{a censura} again
\end{itemize}

\begin{tcolorbox}[breakable,colback=teal!4,colframe=teal!35!black,title={Before you move on}]
What does a censura mean? (\textquotedblleft{}Censorship; criticism\textquotedblright{}.) Say it once more.
\end{tcolorbox}
\section[o editorial]{o editorial — the editorial, the leading article}

\label{lesson:PT-C284-editorial}



\begin{quote}
Before the new one: say the Portuguese for the scandal, then the Portuguese for censorship.
\end{quote}

\subsection*{You'll want to know: o editorial}
\textbf{o editorial} — \textquotedblleft{}the editorial, the leading article\textquotedblright{}.

\begin{quote}
O editorial de hoje critica o governo.
\end{quote}

\begin{quote}
\emph{Today’s editorial criticises the government.}
\end{quote}

As an adjective: \emph{uma decisão editorial}.

\begin{cousinweb}
\emph{editorial} is built on \emph{editor}, a learned borrowing from Latin \emph{ēditor} \textquotedblleft{}one who brings forth, a producer\textquotedblright{} (Lewis \& Short).
\end{cousinweb}

\begin{grammarlens}[title={What you've built}]
One more word for telling, explaining and describing at B1.
\end{grammarlens}

\subsection*{Your turn}
\begin{itemize}
\raggedright
  \item \emph{Say it:} \emph{o editorial}
  \item \emph{Say it:} \emph{o editorial}, once more
  \item \emph{Say it:} a sentence of your own with \emph{o editorial}
  \item \emph{Recall:} say the Portuguese for the scandal, then the Portuguese for censorship, then say \emph{o editorial} again
\end{itemize}

\begin{tcolorbox}[breakable,colback=teal!4,colframe=teal!35!black,title={Before you move on}]
What does o editorial mean? (\textquotedblleft{}The editorial, the leading article\textquotedblright{}.) Say it once more.
\end{tcolorbox}
\section[a crónica]{a crónica — the newspaper piece, a short personal essay; the chronicle}

\label{lesson:PT-C284-cronica}



\begin{quote}
Before the new one: say the Portuguese for censorship, then the Portuguese for the editorial.
\end{quote}

\subsection*{You'll want to know: a crónica}
\textbf{a crónica} — \textquotedblleft{}the newspaper piece, a short personal essay; the chronicle\textquotedblright{}.

\begin{quote}
Leio a crónica dela todos os domingos.
\end{quote}

\begin{quote}
\emph{I read her piece in the paper every Sunday.}
\end{quote}

In Brazil: \emph{crônica}.

\begin{cousinweb}
\emph{crónica} is a learned borrowing from Latin \emph{chronica} \textquotedblleft{}a chronicle\textquotedblright{} (EtymDB), from Greek \emph{khronikós} \textquotedblleft{}of time\textquotedblright{}, built on \emph{khrónos} \textquotedblleft{}time\textquotedblright{}.
\end{cousinweb}

\begin{grammarlens}[title={What you've built}]
One more word for telling, explaining and describing at B1.
\end{grammarlens}

\subsection*{Your turn}
\begin{itemize}
\raggedright
  \item \emph{Say it:} \emph{a crónica}
  \item \emph{Say it:} \emph{a crónica}, once more
  \item \emph{Say it:} a sentence of your own with \emph{a crónica}
  \item \emph{Recall:} say the Portuguese for censorship, then the Portuguese for the editorial, then say \emph{a crónica} again
\end{itemize}

\begin{tcolorbox}[breakable,colback=teal!4,colframe=teal!35!black,title={Before you move on}]
What does a crónica mean? (\textquotedblleft{}The newspaper piece, a short personal essay; the chronicle\textquotedblright{}.) Say it once more.
\end{tcolorbox}
\section[o cronista]{o cronista — the columnist; the chronicler}

\label{lesson:PT-C284-cronista}



\begin{quote}
Before the new one: say the Portuguese for the editorial, then the Portuguese for the newspaper piece.
\end{quote}

\subsection*{You'll want to know: o cronista}
\textbf{o cronista} — \textquotedblleft{}the columnist; the chronicler\textquotedblright{}.

\begin{quote}
O cronista escreve sobre a vida na cidade.
\end{quote}

\begin{quote}
\emph{The columnist writes about city life.}
\end{quote}

The same form for a woman: \emph{a cronista}. In old times, the writer of a kingdom’s history.

\begin{cousinweb}
\emph{cronista} is built on \emph{crónica} \textquotedblleft{}a chronicle, a newspaper piece\textquotedblright{} with \emph{-ista} for the person; \emph{crónica} is a learned borrowing from Latin \emph{chronica}, from Greek \emph{khronikós} \textquotedblleft{}of time\textquotedblright{} (EtymDB).
\end{cousinweb}

\begin{grammarlens}[title={What you've built}]
One more word for telling, explaining and describing at B1.
\end{grammarlens}

\subsection*{Your turn}
\begin{itemize}
\raggedright
  \item \emph{Say it:} \emph{o cronista}
  \item \emph{Say it:} \emph{o cronista}, once more
  \item \emph{Say it:} a sentence of your own with \emph{o cronista}
  \item \emph{Recall:} say the Portuguese for the editorial, then the Portuguese for the newspaper piece, then say \emph{o cronista} again
\end{itemize}

\begin{tcolorbox}[breakable,colback=teal!4,colframe=teal!35!black,title={Before you move on}]
What does o cronista mean? (\textquotedblleft{}The columnist; the chronicler\textquotedblright{}.) Say it once more.
\end{tcolorbox}
\section[(dialogue)]{Revisão dos capítulos 277–284}

\label{lesson:PT-R284-revisao-277-284}



\begin{quote}
Say the Portuguese for frequently and occasionally.
\end{quote}

\subsection*{Your turn}
Say these aloud:

\begin{itemize}
\raggedright
  \item frequently, occasionally, usually, daily, beforehand
  \item the news bulletin, the TV news, the presenter, live, the broadcast
  \item the press, the report, the reporter, the correspondent, the headline
  \item to print, the publication, to make public, the edition, the article
  \item the fact, to report, the account, the written report, to announce in the news
  \item the declaration, the comment, to confirm, the press release, the commentator
  \item the rumour, the murmur, to deny, the denial, the spokesperson
  \item the scandal, censorship, the editorial, the newspaper piece, the columnist
\end{itemize}

\begin{tcolorbox}[breakable,colback=teal!4,colframe=teal!35!black,title={Before you move on}]
Frequently? (\textbf{frequentemente}.) Occasionally? (\textbf{ocasionalmente}.) Usually? (\textbf{habitualmente}.) Daily? (\textbf{diariamente}.) Beforehand? (\textbf{previamente}.) The news bulletin? (\textbf{o noticiário}.) The TV news? (\textbf{o telejornal}.) The presenter? (\textbf{o apresentador}.) Live? (\textbf{em direto}.) The broadcast? (\textbf{a transmissão}.) The press? (\textbf{a imprensa}.) The report? (\textbf{a reportagem}.) The reporter? (\textbf{o repórter}.) The correspondent? (\textbf{o correspondente}.) The headline? (\textbf{a manchete}.) To print? (\textbf{imprimir}.) The publication? (\textbf{a publicação}.) To make public? (\textbf{divulgar}.) The edition? (\textbf{a edição}.) The article? (\textbf{o artigo}.) The fact? (\textbf{o facto}.) To report? (\textbf{relatar}.) The account? (\textbf{o relato}.) The written report? (\textbf{o relatório}.) To announce in the news? (\textbf{noticiar}.) The declaration? (\textbf{a declaração}.) The comment? (\textbf{o comentário}.) To confirm? (\textbf{confirmar}.) The press release? (\textbf{o comunicado}.) The commentator? (\textbf{o comentador}.) The rumour? (\textbf{o boato}.) The murmur? (\textbf{o rumor}.) To deny? (\textbf{desmentir}.) The denial? (\textbf{o desmentido}.) The spokesperson? (\textbf{o porta-voz}.) The scandal? (\textbf{o escândalo}.) Censorship? (\textbf{a censura}.) The editorial? (\textbf{o editorial}.) The newspaper piece? (\textbf{a crónica}.) The columnist? (\textbf{o cronista}.)
\end{tcolorbox}
